\section{Scope, classification targets, and open boundaries}
\label{sec:publication-scope}

The global classification theorem is componentwise because the certified reversible core has five connected semantic components.  The component tag is part of the invariant.  Omitting it would compare objects for which no reversible semantic bridge has been proved.

The lane theorems have a different scope.  The finite profiled Universal Lane Transport Theorem is exact inside a declared additive, idempotent-complete setting with complete split packet systems.  Its realization, kernel, multilinear, fold, and descent clauses apply only under their additional hypotheses.  Trivial address holonomy is required only for one fixed global address set.  With a coherent finite address cocycle, complete matrices instead descend as sections of a matrix local system, and monodromy orbits give the finest ordinary global projector coarsening.  The word ``universal'' refers to this formal lane calculus.  It does not claim that every object has a finite lane presentation or that address count equals dimension, rank, entropy, information, or computational complexity.

The semantic quotient theorem is relative in another direction.  It fixes a carrier doctrine, continuation grammar, observation category, and query family.  The resulting quotient is universal only for that semantic slice.  A different observation category can identify different linear combinations of residual states.  A different carrier can change the legal aggregation law.  A different continuation class can expose provenance that was previously invisible.  There is therefore no observer-independent minimum hidden behind the notation.

The admissibility theory records a third boundary.  A positive defect in a quadratic, integral, finite, or neutral class need not survive enlargement to a convex, rational, completed, or nonneutral class.  Conversely, zero optimized defect does not imply that a completion exists.  Attainment and conditioning remain part of the defect profile.

\begin{landscape}
\begin{table}[htbp]
\centering
\scriptsize
\caption{Certified classification components and exact boundaries.}
\label{tab:block14-scope}
\begin{tabularx}{\linewidth}{@{}L{0.15\linewidth}L{0.23\linewidth}L{0.30\linewidth}Y@{}}
\toprule
Component & Semantic domain & Complete descended coordinates & Boundary \\
\midrule
Digital--regular--linear--Mahler & Fixed-prime regular sequences with rational-mask and periodic subloci & Strict shortlex state, Cartier/mask data where defined, periodic linear data where defined & No all-radix or unrestricted functional/Newton-Mahler equivalence \\
Continuous local & $C(\mathbb Z_p,\mathbb Z_p)$ & Canonical Mahler coefficients and complete clopen branch packets & No universal local-analytic normal form or optimal mixed-radix tree \\
Marked Fuss & $U=z+\alpha U^q$ & Marked parameters $(q,\alpha)$, algebraic branch, recurrence, tree support and orbit data & Complete only on the marked nonlinear--algebraic--holonomic family \\
Euler--multiplicative & Fixed Euler basis and prime-uniform multiplicative locus & Euler exponent, divisor ghost, common local factor, multiplicative function & No arbitrary product-grammar or additive--multiplicative equivalence \\
Finite Galois & One finite semisimple Galois extension & Subgroup, fixed field, Reynolds projector, quotient and Mackey data & No modular/wild or nonneutral tensor reconstruction \\
\bottomrule
\end{tabularx}
\end{table}
\end{landscape}


\subsection{Directional satellites}

Several useful constructions remain adjacent to the reversible core without becoming equivalences:
\begin{itemize}
\item reduction from integral section lattices to finite residue systems, with Smith--Tor ghosts and missing saturated coordinates;
\item passage from a selected algebraic branch to a differential annihilator, which can forget the branch and its minimal polynomial;
\item functional-Mahler masks and Newton--Mahler interpolation, which use different variables and require additional lifting and interpolation data;
\item projective Galois descent, whose lift to an ordinary vector realization is governed by a relative Brauer class;
\item faithful additive or Hilbert-space realizations that do not supply the exact tensor fiber functor required for Tannakian reconstruction;
\item incomplete lane selections and noninjective folds, which retain explicit residuals or collision kernels;
\item Boolean or max-valued summaries of signed additive computation, where cancellation obstructs carrier replacement;
\item Set-valued future quotients and weighted-linear residual spaces, whose difference is measured by a relation module;
\item geometric or provenance compactions that are valid only for the declared continuation class and downstream task.
\end{itemize}
These constructions remain in the directed transformation category with their exact loss, completion, or obstruction records.

\subsection{Open mathematical obligations}

The remaining foundational obligations are more sharply localized than in earlier editions.
\begin{enumerate}
\item Extend the certified atlases across radices, coefficient rings, semirings, and mixed local--global completions without losing normalization or effectivity.
\item Develop modular and wild analogues of finite semisimple Galois projector theory when Reynolds averaging fails.
\item Replace marked model-family results in the holonomic, algebraic, and nonlinear strata by broader native minimality and collision invariants.
\item Construct nonneutral and integral tensor-descent variants with their gerbe or torsor data retained explicitly.
\item Prove finite-data recognition theorems only on declared subgrammars where identifiability and certification are available.
\item Extend finite split and finite holonomy-twisted lane transport to infinite, completed, nonsplit, noninvertible, or genuinely nonadditive resolutions only after specifying the required convergence, exactness, and coherence doctrines.
\item Characterize which query universes are rich enough that every continuation-stable congruence is observationally definable, and determine when the definable task family is closed under joins.
\item Extend the module exact-sequence calculus to exact, regular, homological, enriched, and higher-categorical settings without silently importing additive arguments.
\item Develop effective algorithms for joint semantic completion, relation-module reduction, and task-relative bridge loss beyond finite matrices and finite automata.
\item Construct a fully coherent doctrine-indexed comparison object only after the required horizontal, vertical, and higher comparison cells have been verified.
\end{enumerate}

\begin{boundarybox}
\textbf{No global promotion.}
Completion of the finite atlas, lane theorems, semantic quotient calculus, or class-relative defect theory does not promote absent bridges, incompatible carriers, noninjective folds, unresolved enrichments, or broad-stratum normal forms to equivalences.  Every claimed transport remains relative to its grammar, profile, semantic slice, admissible realization class, and proof-status record.  The detailed obligation and claim ledgers are retained in \cref{app:block14-open-program,app:block14-claim-ledger}.
\end{boundarybox}
